\section{Source File Size}
Above a certain number of lines, source code files become unwieldy, but it is difficult to specify a concrete value for the maximum size of a Java source file. First, generated source files tend to have lots of lines because it is often much easier to use boiler-plate code here than to generate methods and call them instead. And in particular when dynamically generated code is used\footnote{Code that is regenerated for each build.}, it would not even hurt the DRY principle: a change to the repeated code can be performed on one single place, the generator.

Second, classes that have lots of attributes usually end up with lots of lines, too: the formatting style and coding guidelines proposed in this document require a least 5~lines for each attribute, 6~lines for each getter method and 11~lines for each setter method, always including all comments and blank lines. This is about 22~lines for each attribute.  One hundred attributes seems to be a lot, but I have seen more than one real life application where database tables had more than that number of columns that each needed to be reflected as an attribute of the class that should reflect the table. I do not think that this is a good design anyway, but if your new Java class has to reflect the legacy data model, it has to have that number of attributes, bringing you to about 2200~lines of code without any methods other than the already mentioned getters and setters.

I found a limit of about 2000~lines to be feasible; if a class becomes larger than this number of lines, you should consider to move some of the code into other classes. This could be a utility class\index{utility class}, a parent class (probably abstract), or the functionality would be better placed with another entity in general. Also externalising some functionality by implementing the command pattern\autocite{Gamma:DesignPatterns} might be an idea. 

\subsubsection{Principles}
\begin{enumerate}[label=P\arabic*.]
\item{A Java source file should not exceed the limit of 2000~lines. As shorter it is, as better.}
\item{Generated source files may have significant more lines.}
\item{Other formatting rules may not be sacrificed only to reduced the number of lines of a source file.}
\end{enumerate}

%==============================================================================
% End of file!!
%==============================================================================
